\section{Conclusion}
\label{sec:conclusion}

\tothink{full conclusion prose; outline below.}

\noindent We presented a three-axis methodology for LLM-as-judge eval-suite
scoring. The methodology partitions the (scorer-class $\times$
task-shape) grid into cells whose primary reliability axis is
\emph{calibration} (LLM classifier, pairwise rubric), \emph{agreement}
(LLM rubric, trajectory rubric on subjective shapes), or \emph{neither
applies} (heuristic scorers, where dataset quality is upstream of the
scorer). Every LLM-judge cell additionally satisfies a \emph{bias-bound
precondition} on position, verbosity, and self-preference biases. The
three-axis frame is structurally distinct from the prevailing
single-axis framing in the eval literature, which collapses
operationally-distinguishable failure modes that need different fixes.

\noindent The paper's central methodological claim is that
\emph{eval-suite reliability is a three-axis problem with non-collapsible
axes}, structurally distinct from the tier-cost-ratio framing of the
companion model-advisor paper. The paper's central empirical claim,
following Kadavath et al.\ \cite{kadavath2022language} and reproduced
on judge tasks specifically, is that \emph{RLHF-tuned LLM judges are
systematically worse-calibrated than their pretrained ancestors at
the same scale}, with operational implications for any team running
forced-tool-use LLM judging on instruction-tuned models. Probability-weighted
scoring \cite{liu2023geval} and temperature scaling
\cite{guo2017calibration} are the cheap fixes; the eval-suite ships
both as v1 defaults.

\paragraph{Acknowledgements.}
\tothink{contributors, operators, and the open-source artefacts. Same
shape as the model-advisor and retriever papers.}
